% Definitional completion R-032; provenance in
% spec/research/edition-2026-c4-pa2-definition.md.
\begin{hbdefinition}{def:p02-pa2}{The finite full second-order arithmetic sentence}
For the full-validity reduction below, use the signature
\(\{0,S,+,\times\}\), where \(0\) is a constant, \(S\) is unary,
and \(+\) and \(\times\) are binary function symbols; equality has its
logical interpretation.  The notation \(\mathsf{PA}_2\) denotes the
conjunction of the following seven closed sentences:
\[
\begin{gathered}
  \forall x\,(Sx\ne0),
  \qquad
  \forall x\,\forall y\,(Sx=Sy\to x=y),\\
  \forall X\,
    \bigl[(X(0)\land\forall x\,(X(x)\to X(Sx)))
           \to\forall x\,X(x)\bigr],\\
  \forall x\,(x+0=x),
  \qquad
  \forall x\,\forall y\,(x+S(y)=S(x+y)),\\
  \forall x\,(x\times0=0),
  \qquad
  \forall x\,\forall y\,
    (x\times S(y)=(x\times y)+x).
\end{gathered}
\]
The variable \(X\) is unary.  Under the full semantics of
\cref{sec:second-order-logic-02}, its quantifier ranges over every
subset of the nonempty object domain; object quantifiers range over that
domain, and the function symbols denote total functions of their declared
arities.  The induction conjunct is one second-order sentence, rather than
a first-order induction scheme.

The first three conjuncts are the successor and induction axioms displayed
in \cref{sec:second-order-logic-03}; the final four make explicit the
addition and multiplication recursion axioms used in its expanded-signature
conclusion.  Thus \(\mathsf{PA}_2\to\sigma\), for a first-order sentence
\(\sigma\) in this signature, denotes a single second-order sentence with
a fixed finite antecedent.  The notation here includes the four arithmetic
conjuncts; it is not an abbreviation for the successor-only axiom set.
\index{second-order arithmetic}
\index{Peano axioms!full second-order}
\end{hbdefinition}
